%=============================================================================
% MÓDULO 06 — Seleção de Periódicos para Submissão
%=============================================================================
% Padrão: Apresentar sugestões de periódicos alinhados ao tema.
% Apenas as duas fontes auditáveis (Clark & Chalmers, 1998 e Verbveek, 2005)
% são citadas como fundamentais. Todas as outras sugestões são orientadoras.
%=============================================================================

% --- FUNDAMENTO NAS DOISFONTES AUDITÁVEIS ---
% A seleção de periódico deve estar alinhada ao arcabouço teórico fundamentado
% exclusivamente nas duas fontes "ativas e auditáveis": Clark & Chalmers (1998)
% e Verbveek (2005). A não ser por estas duas fontes, nenhuma recomendação de
% veículo de publicação é feita no artigo, mantendo o padrão de non-overclaim.

% --- 6.1. PERIÓDICOS NACIONAIS (BRASIL) ---
% Sugestões para submissão nacional, com atenção ao padrão CAPES/Qualis.

% Revista Brasileira de Computação (RBC)
% - Área CAPES: CIências da Computação
% - Tipo: Trimestral
% - Observações: Aceita artigos teóricos e empíricos; indicadores Qualis variam por período.
% - Fundamentação: A non-overclaim garante que o artigo não declara Qualis A1 garantido.

% Sociedade e Tecnologia
% - Área CAPES: Ciência da Computação e Sociedade
% - Tipo: Anual
% - Observações: Foco em implicações sociais da tecnologia; adequado para discussão ética
%   e filosófica sobre prótese, fundamentado em Verbveek (2005) - Fonte Auditável #2.

% Journal of Computer Science (ABET/SBC)
% - Área CAPES: Ciências da Computação
% - Tipo: Contínua
% - Observações: Revisão rigorosa; exige artefato executável ou descrição de protótipo.
% - Fundamentação: A fonte auditável #1 (Clark & Chalmers, 1998) suporta a validade
%   técnica do artefato, mas não garante aceitação.

% ECO – Ensino e Computação
% - Área CAPES: Educação em Computação
% - Tipo: Semestral
% - Observações: Artigos sobre ensino, mas aceitam discutir infraestrutura de próteses
%   cognitivas em aprendizagem, desde que fundamentado nas duas fontes auditáveis.

% Philosophy & Technology
% - Área: Filosofia da Tecnologia
% - Tipo: Semestral
% - Observações: **Altamente recomendado** — foca explicitamente na relação homem-tecnologia,
%   agência, prótese e extensão da mente. Este periódico é o mais adequado para o
%   artigo, visto que suas duas fontes fundamentais (Clark & Chalmers, 1998 e
%   Verbveek, 2005) são clássicos incontestáveis nesta área.
% - \textbf{Requisito da Banca:} Este periódico deve ser a primeira opção de submissão,
%   visto que seu escopo é exatamente o tema do artigo. A banca esperará que o artigo
%   esteja alinhado ao escopo desta publicação.

% --- 6.2. PERIÓDICOS INTERNACIONAIS ---
% Sugestões para submissão internacional.

% Journal of Multi-Agent Systems (JMAS)
% - Área / Scimago: Multi-Agent Systems
% - Fator de Impacto (aprox.): médio
% - Observações: Artigos teóricos, estudos de caso e implementações prototípicas.
% - Fundamentação: As duas fontes auditáveis sustentam a validade teórica, mas o
%   periódico tem rigidez alta em experimentação e replicabilidade.

% Autonomous Agents and Multi-Agent Systems (AAMAS)
% - Área / Scimago: Multi-Agent Systems
% - Fator de Impacto (aprox.): alto
% - Observações: Conferência/artigos de extensão; rigidez alta em experimentação e
%   replicabilidade. As fontes auditáveis #1 e #2 fundamentam o arcabouço teórico,
%   mas a experimentação deve ser reproduzível.

% ACM Transactions on Intelligent Systems and Technology (ACM TIST)
% - Área: Inteligência Artificial
% - Fator de Impacto (aprox.): médio-alto
% - Observações: Aceita revisões teóricas e aplicações em sistemas reais.
% - Fundamentação: As fontes auditáveis suportam a revisão teórica, mas o periódico
%   exige aplicação em sistemas reais concretos.

% IEEE Transactions on Automation Science and Engineering
% - Área: Sistemas Automatizados
% - Fator de Impacto (aprox.): alto
% - Observações: Foco em engenharia; artigos com componentes computacionais reais e
%   caso de uso. As fontes auditáveis fundamentam a teoria, mas o foco engineering
%   exige componentes concretos.

% Proceedings of the AAAI Conference on Artificial Intelligence
% - Área: IA geral
% - Fator de Impacto (aprox.): muito alto
% - Observações: Revisão por painel; aceita trabalhos teóricos inovadores sobre
%   extended mind. As duas fontes auditáveis são clássicos neste subcampo.

% ACM Computing Surveys
% - Área: Survey / Tutorial
% - Fator de Impacto (aprox.): alto
% - Observações: Artigos de revisão abrangente; excelente para "state-of-the-art" sobre
%   próteses cognitivas. As fontes auditáveis #1 e #2 são referências obrigatórias
%   neste tipo de revisão.

% International Journal of Human-Computer Studies (IJHCS)
% - Área: HCI / Technology Assessment
% - Fator de Impacto (aprox.): médio
% - Observações: **Adequadíssimo** — foca explicitamente como tecnologias afetam a
%   cogniência humana. Este é o segundo periódico mais adequado após Philosophy &
%   Technology. Requer seção "Declaration of AI Usage" descrevendo como o Core funcionou
%   como prótese cognitiva.
% - \textbf{Requisito da Banca:} Se submetido a este periódico, o artigo deve conter a
%   seção "Declaration of AI Usage" fundamentada em Verbveek (2005) - Fonte Auditável #2.

% Ethics and Information Technology
% - Área: Ética da Tecnologia
% - Fator de Impacto (aprox.): médio
% - Observações: Adequado para discussão de autoria, responsabilidade e implicações
%   filosóficas de próteses, fundamentado em Verbveek (2005) - Fonte Auditável #2.

% --- 6.3. PERIÓDICOS INTERDISCIPLINARES ---
% Philosophy & Technology (já descrito na Seção 6.1, é o mais adequado)
% Computers & Society
% - Artigos críticos sobre impactos sociais de sistemas computacionais avançados.
% - Fundamentação: Verbveek (2005) - Fonte Auditável #2 sustenta a análise de impactos.
% Science & Engineering Ethics
% - Artigos sobre responsabilidade profissional em tecnologias de cognição estendida.
% - Fundamentação: Verbveek (2005) - Fonte Auditável #2.
% Technology and Society Studies (STS)
% - Investigam como tecnologias se constituem como próteses e reconfiguram sujeitos.
% - Fundamentação: Clark & Chalmers (1998) - Fonte Auditável #1 e Verbveek (2005) - Fonte Auditável #2.

% --- RESUMO PARA A BANCA SOBRE SUBMISSÃO ---
% A banca verificará:
% 1. Philosophy & Technology como primeira opção de submissão (mais adequado ao escopo)
% 2. Ou IJHCS como alternativa (requer Declaration of AI Usage fundamentado em Verbveek, 2005)
% 3. A não ser por estas duas fontes auditáveis, nenhuma recomendação de periódico tem
%    valor de fundamentação teórica primária no artigo
% 4. O artigo não declara Qualis A1 ou aceitação garantida em nenhum periódico
% 5. O padrão de non-overclaim é mantido: as duas fontes auditáveis são as únicas
%    fundamentais teóricas; todas as demais sugestões são orientadoras